% Section 10 (writer: steep-a). Prose summary of Paper 2 (claudeSFI2026steeper); no theorem environments, no new claims.
% Every number matches ../paper-steeper/CLAIMS.md. Paper 2's letters: programs pi, lengths ell(pi); per-sample rate written lambda/n.
% Sources (through Paper 2): research/tracks/steeper/notes-final.md §§1-4.
\section{Steeper simplicity priors (summary of the companion paper)}\label{sec:steep}

H\"anni's notes propose to penalise ``the length of the program plus the negative log likelihood \dots{} but with a higher coefficient on the length of the program'', in function induction where ``each input gets its own private randomness'', so that the selected hypothesis is a simple model plus noise rather than a model plus memorised errors of the teacher \citep{hanni2026notes,hanni2026lesswrong}. The companion paper \citep{claudeSFI2026steeper} studies this proposal. It works with classes of hypotheses indexed by programs $\pi$ of length $\ell(\pi)$, the objective $\lambda\ell(\pi)-\log P_\pi(D)$ with $\lambda\ge1$, and the power posterior $\propto2^{-\lambda\ell(\pi)}P_\pi(D)$. With $\lambda$ proportional to the sample size $n$ these are batch objectives at each $n$, not one Bayesian predictor. This section summarises the results; all of them are proved there.

\paragraph{The absorption objection.} A standard point, which H\"anni's note already makes for sequence prediction, is that under a steeper penalty a hypothesis that is itself a universal mixture absorbs everything. The companion paper's \emph{absorption theorem} makes this exact: in any class that contains its own mixture $U$ as a member of length $\ell_U$ (SI; SFI with shared randomness, $U=\MF$; the weak class, $U=\Mw$; \SFIhist{}, $U=\MH$, which is universal by \cref{cor:hist:H4}), every $\lambda$-MAP has length at most $(\lambda\ell_U+O(1))/(\lambda-1)$ on every data set and fits almost as well as $U$, so a deterministic teacher's errors are reproduced. In the class of product hypotheses (private randomness per input, H\"anni's x-only samplers) no member is universal: every fixed product hypothesis loses $2^d-1$ bits on some parity on $\{0,1\}^d$ that has a program of length $d+2\log d+O(1)$, and the universal conditional semimeasure never learns. So steeper SFI is defined over the product class, whose $\lambda\to1$ limit is $\Mprod$, not $\MF$. This restriction has a price. At $\lambda=1$ it is the gap between $\Mprod$ and $\MF$, unbounded but at most $K(\dom D)+O(1)$ bits (\cref{prop:versions:shared,rem:versions:sharedsize}). At $\lambda>1$, on labels that share a latent fair coin between pairs of inputs, the product class pays $\min(\lambda,2)$ bits per coin against $1$ bit for $\MF$. The protection ends where hypotheses see earlier \emph{labelled} pairs: the in-context Bayes mixture is then one product hypothesis again. The paper presents the restriction to private randomness (per input; on data sets the same as per query, \cref{prop:learn:perquery}) as a modelling choice: the model is the part of the description reused across inputs.

\paragraph{The teacher's mistakes.} H\"anni's example has a $100$-bit simple model and systematic errors on a $2^{-10}$ fraction of the inputs whose set takes $10{,}000$ bits. The companion paper realises it rigorously with a random parity on $\{0,1\}^{100}$ and an error rule on the first $20$ input bits ($11{,}710$ bits), quantifying over \emph{all} product hypotheses by list decoding. For error rules that are incompressible relative to the model, the \emph{teacher's-mistakes theorem} gives a window of $\lambda/n$ from $2^{-20}\approx9.54\cdot10^{-7}$ up to $8.06\cdot10^{-3}$ (machine constant set to $0$, an illustration) or $3.66\cdot10^{-3}$ (machine constant $30$). Inside it every near-optimal hypothesis knows the model and does not reproduce the errors; at $\lambda/n=10^{-4}$, with the machine constant set to $0$, it gives the parity's label average probability at least $0.983$, holds at most $37.7$ of the $11{,}710$ bits of the error rule, and gives the erroneous label average probability at most $0.199$ on the error inputs. A few per cent below $2^{-20}$ and further down, near-optimal hypotheses reproduce the errors. The theorem uses a list-decodable model class, the product form of the error set, uniform inputs, an error density below $0.056$, and a small randomness deficiency of the error rule. Plain Bayes ($\Mprod$, $\MF$), SI, and the $\lambda$-MAP under every schedule with $\lambda/n\to0$ learn the errors, so $\lambda$ must grow linearly in $n$, and the learner is then deliberately not consistent for the teacher's label distribution. That a bounded $\lambda$ overfits and consistency needs $\lambda$ to grow is known for $\lambda$-weighted two-part MDL in noisy classification \citep{grunwald2007suboptimal,zhu2025quantifying,zhu2026overfitting}, where the target is the data distribution rather than a norm that differs from it.

\paragraph{Limits.} The trick sorts components by rate, not by validity. Density and description length do not suffice: with H\"anni's density and an error set made of a cheap subcube of $128$ points plus $896$ random points (about $10{,}400$ bits), no hypothesis that fits no better than the good model is near-optimal for $\lambda/n<1.6\cdot10^{-5}$, about $31\%$ of the window's logarithmic width ($25\%$ with machine constant $30$). This refutes the extension of H\"anni's example, which he called ``somewhat cherry-picked'', to every error rule with his density and description length; whether structured error rules leave a window is open. A cheap, frequent error rule defeats the trick for all input lengths above a threshold linear in the machine constant (rate blindness); $\lambda$ cannot be chosen from imitation data, and the window depends on the input law. Out of distribution, $\lambda$ proportional to $n$ removes the constant lower bound of \cref{prop:ood:hanni}, which becomes $2^{-\lambda(K(Q,b)+O(1))}$, but only by sharpening the preference for the shortest hypotheses that fit the training data; no guarantee about unseen inputs follows.

\paragraph{The model--randomness split.} H\"anni also asked for ``a preference for a simple sampler over a model where the sampled bits are filled in'' and noted that Solomonoff induction has none, ``both get similar posterior'' \citep{hanni2026request}. At $\lambda=1$ this is confirmed with explicit factors for every computable sampler $\nu$ with a program of length $k$: in SFI the posterior odds of the deterministic fill-ins against the sampler lie in $[2^{-2\log(t_n+1)-O(1)},2^{k+\delta^*}]$, with $t_n$ doubly logarithmic in the input length and $\delta^*$ a Ville deficiency with $\Pr[\delta^*\ge j]\le2^{-j}$, whatever the amount of sampled randomness; this covers non-dyadic samplers such as the Bern$(1/3)$ threshold hypotheses of Eisenstat's example (\cref{def:ex:eisenstat}). SI behaves the same up to $2^{O(k+K(n))}$. For every fixed $\lambda>1$ in every countable product class, the sampler wins: every $\lambda$-MAP has length at most $k+(k+\delta^*)/(\lambda-1)$, the fill-ins' posterior is at most $2^{\lambda k+\lambda\delta^*-(\lambda-1)(-\log\nu(D_n))}$, and with infinitely many distinct inputs the $\lambda$-MAPs are eventually equivalent to $\nu$ on the inputs seen (summable Hellinger distances, not equality). With shared randomness, or in SI, the universal mixture wins instead: no sampler with $(\lambda-1)k>\lambda\ell_U+O(1)$ is ever selected. In the language of algorithmic statistics, for a noisy random parity on the full truth table the set (shared-randomness) sophistication is $O(\log d)$, so the parity's parameter counts as noise, while the product sophistication is $d\pm O(\log d)$ for every flip rate below $1/4$: the product structure function has the kink at the simple model that H\"anni anticipated, where with shared randomness ``1 bit pays for 1 bit''.
